\textbf{Lower image likelihood loss does not ensure nonempty conditional writing.}
\textbf{a,b,} Mean source-mask negative log likelihood and false-blank counts
for the same complete supplied-history book groups, comparing the head-only
control with learned continuation context. NLL chart/workbook values round to
three decimals; full precision remains in the bound source data. Counts are
exact. False blanks exclude genuinely empty sources, counted in the table.
Every group retains all positions; development has 1,161 and each training
layout 4,586. \textbf{c,} Complete-image source preference among false blanks,
using the fixed 0.001-nat margin. The observed counts 4/93, 26/307 and 27/289
miss the prospective minimum counts 47, 154 and 145. Across all seven groups,
57/711 false blanks prefer the source; the zero-prefix upper bound supports
an empty global mode at 108 false blanks, including 102 book positions.
The numerical inequality is subject to finite precision and is not a formal
certificate. All saved physical argmax choices reproduce.
\textbf{d,} The first six of 21 development-book certified empty-mode failures,
ordered by original case and zero-based position. Each triplet retains source,
head-only and learned-context binary pixels, expanded threefold by nearest
neighbour without repair. All source samples contain ink and all shown context
outputs are empty. The complete probability record retains all eligible cases.
Original review grades remain frozen despite 28 contradictory blank flags,
including four candidate and two control credited empty images. One model
reviewer had prior cohort familiarity; independent human adjudication is absent.
The failed physical prerequisite prevents autonomous evaluation. Neither lower
likelihood loss nor this diagnostic establishes complete language, a compact
foundation or matched-quality efficiency.
